% Chapter 1: Core definitions
\chapter{Definitions}

\section{Operator Convexity and Concavity}

Throughout, $\mathcal{B}$ denotes a unital C$^*$-algebra equipped with a compatible
partial order (i.e.\ a \emph{ordered} C$^*$-algebra), and $I \subseteq \mathbb{R}$
denotes a subset of the real line.
For a continuous function $f : \mathbb{R} \to \mathbb{R}$, we write $f(a)$ for the
element of $\mathcal{B}$ obtained by applying the continuous functional calculus to
a self-adjoint element $a \in \mathcal{B}$.

\begin{definition}[Operator convex function]
  \label{def:operator_convex}
  \lean{OperatorConvexOn}
  \leanok
  A function $f : \mathbb{R} \to \mathbb{R}$ is \emph{operator convex} on $I$ if, for
  every ordered unital C$^*$-algebra $\mathcal{B}$, the map $a \mapsto f(a)$ is convex
  on the set
  \[
    \{ a \in \mathcal{B} \mid a \text{ is self-adjoint},\ \sigma(a) \subseteq I \},
  \]
  where $\sigma(a)$ denotes the spectrum of $a$ over $\mathbb{R}$.
\end{definition}

\begin{definition}[Operator concave function]
  \label{def:operator_concave}
  \lean{OperatorConcaveOn}
  \leanok
  A function $f : \mathbb{R} \to \mathbb{R}$ is \emph{operator concave} on $I$ if
  $-f$ is operator convex on $I$, equivalently if $a \mapsto f(a)$ is concave on the
  set of self-adjoint elements of every ordered unital C$^*$-algebra with spectrum
  contained in $I$.
\end{definition}

\section{Generalized Perspective Function}

\begin{definition}[Generalized perspective function]
  \label{def:gen_perspective}
  \lean{GenPerspective}
  \leanok
  Let $\mathcal{B}$ be an ordered unital C$^*$-algebra and let $f, g : \mathbb{R} \to \mathbb{R}$.
  The \emph{generalized perspective function} $(f \triangle g) : \mathcal{B} \times \mathcal{B} \to \mathcal{B}$
  is defined by
  \[
    (f \triangle g)(L, R)
    \;:=\;
    g(R)^{1/2}\, f\!\left(g(R)^{-1/2}\, L\, g(R)^{-1/2}\right) g(R)^{1/2}.
  \]
\end{definition}

\section{H\texorpdfstring{$^*$}{*}-Algebra}

\begin{definition}[H$^*$-algebra]
  \label{def:hstar_algebra}
  \lean{HStarAlgebra}
  \leanok
  An \emph{H$^*$-algebra} over $\mathbb{k} \in \{\mathbb{R}, \mathbb{C}\}$ is a
  Hilbert space $H$ over $\mathbb{k}$ equipped with the structure of a normed
  $\mathbb{k}$-algebra with involution $\star$, satisfying the compatibility axioms:
  \[
    \langle a x,\, y \rangle = \langle x,\, a^* y \rangle
    \qquad \text{and} \qquad
    \langle x a,\, y \rangle = \langle x,\, y a^* \rangle
    \qquad \forall\, a, x, y \in H.
  \]
  This notion was introduced by Ambrose~\cite{Ambrose1945}.
\end{definition}

The canonical example is $H = M_n(\mathbb{C})$ with the Frobenius inner product
$\langle A, B \rangle = \operatorname{tr}(A^* B)$.

\begin{definition}[Left and right multiplication operators]
  \label{def:lmul}
  \lean{Lmul}
  \leanok
  For an H$^*$-algebra $H$ and an element $a \in H$, the \emph{left multiplication
  operator} $L_a : H \to H$ is the bounded linear map $x \mapsto a x$.

  \lean{Rmul}
  The \emph{right multiplication operator} $R_a : H \to H$ is the bounded linear map
  $x \mapsto x a$.

  Both $a \mapsto L_a$ and $a \mapsto R_a$ are continuous as maps $H \to (H \to_L H)$.
  Moreover $L_a$ and $R_b$ commute for all $a, b \in H$:
  $L_a \circ R_b = R_b \circ L_a$.
\end{definition}

The maps $a \mapsto L_a$ and $a \mapsto R_a$ extend to star algebra homomorphisms:
$a \mapsto L_a$ is a $*$-homomorphism $H \to (H \to_L H)$, and
$a \mapsto R_a$ is a $*$-homomorphism $H \to (H \to_L H)^{\mathrm{op}}$
(an anti-homomorphism on $H$, since $R_{ab} = R_b \circ R_a$).

\begin{proposition}[CFC commutativity of $L$ and $R$]
  \label{prop:lmul_map_cfc}
  \lean{Lmul_map_cfc, Rmul_map_cfc}
  \leanok
  \uses{def:lmul, def:hstar_algebra}
  Let $H$ be an H$^*$-algebra, $a \in H$ self-adjoint, and $f : \mathbb{R} \to \mathbb{R}$
  continuous on $\sigma(a)$. Then:
  \[
    L_{f(a)} = f(L_a), \qquad R_{f(a)} = f(R_a),
  \]
  where $f(L_a)$ and $f(R_a)$ denote the continuous functional calculus applied to
  the operators $L_a, R_a \in (H \to_L H)$.

  As a consequence, for any $r \in \mathbb{R}$:
  \[
    (L_a)^r\, x = a^r x, \qquad (R_a)^r\, x = x\, a^r,
  \]
  whenever $a \geq 0$ (for $r \geq 0$) or $a$ is strictly positive (for all $r$).
\end{proposition}
